%---------------------------------------------------------------------
\begin{glyph}{Member (員)}{member}\label{glyph:member}
Member.
\end{glyph}
\gindex{Member}{員}\strindex{10}{員}

\begin{comps}
\comprtwocone&口 + 貝\\\hline
\comprthreecone&Mouth + Shell\\\hline
\end{comps}

\begin{remglyph}
\hooglig{Member} of the \remcom{vampire} club collecting \remcom{shells}.\\\hline
\end{remglyph}

\begin{prons}
\pronsrtwocone&\textbf{イン (IN)}&---\\\hline
\pronsrthreecone&---&---\\\hline
\end{prons}

\begin{remprons}
\hooglig{A member} of the \comon{inklings}.\\\hline
\end{remprons}

%{\middel}
%&\mbox{()}\\\hline
\begin{somme}
全員&complete + member = \textit{all members; the whole staff}&ゼン{\middel}イン \mbox{(ZEN{\middel}IN)}\\\hline
店員&shop + member = \textit{clerk; shop assistant}&テン{\middel}イン \mbox{(TEN{\middel}IN)}\\\hline
社員&company + member = \textit{employee; member}&シャ{\middel}イン \mbox{(SHA{\middel}IN)}\\\hline
工員&craft + member = \textit{factory worker}&コウ{\middel}イン \mbox{(K\={O}{\middel}IN)}\\\hline
動員&move + member = \textit{mobilisation}&ドウ{\middel}イン \mbox{(D\={O}{\middel}IN)}\\\hline
教員&teach + member = \textit{teacher; instructor}&キョウ{\middel}イン \mbox{(KY\={O}{\middel}IN)}\\\hline
行員&go + member = \textit{bank clerk}&コウ{\middel}イン \mbox{(K\={O}{\middel}IN)}\\\hline
定員&fixed + member = \textit{fixed number; capacity}&テイ{\middel}イン \mbox{(TEI{\middel}IN)}\\\hline
要員&essential + member = \textit{necessary personnel}&ヨウ{\middel}イン \mbox{(Y\={O}{\middel}IN)}\\\hline
研究員&research + member = \textit{lab worker}&ケン{\middel}キュウ{\middel}イン \mbox{(KEN{\middel}KY\={U}{\middel}IN)}\\\hline
団員&group + member = \textit{group member}&ダン{\middel}イン \mbox{(DAN{\middel}IN)}\\\hline
欠員&lack + member = \textit{vacancy}&ケツ{\middel}イン \mbox{(KETSU{\middel}IN)}\\\hline
係員&connection + member = \textit{person in charge}&かかり{\middel}イン \mbox{(kakari{\middel}IN)}\\\hline
員数&member + number = \textit{number of members}&インズウ \mbox{(IN{\middel}Z\={U})}\\\hline
組員&association + member = \textit{gang member}&くみ{\middel}イン \mbox{(kumi{\middel}IN)}\\\hline
通信員&correspondence + member = \textit{correspondent; reporter}&ツウ{\middel}シン{\middel}イン \mbox{(TS\={U}{\middel}SHIN{\middel}IN)}\\\hline
\end{somme}

\end{multicols}
\begin{decomp}{6}
日本&の&会社員&は&良く&働く。\\\hline
ニッポン&の&カイシャイン&は&とく&はたらく\\\hline
Japan&の&company employee&は&skilfully&to work\\\hline
\multicolumn{6}{|r|}{\textit{Japanese office workers work very hard.}}\\\hline
\end{decomp}
\pagebreak
%%--------------------------------------------------------------------
%1--------------------------------------------------------------------
\begin{multicols}{2}
\begin{glyph}{Cargo (荷)}{cargo}\label{glyph:cargo}
Cargo.  Burden.
\end{glyph}
\gindex{Cargo}{荷}\strindex{10}{荷}

\begin{comps}
\comprtwocone&艹 + 何\\\hline
\comprthreecone&Grass + What (p. \pageref{glyph:what})\\\hline
\end{comps}

\begin{remglyph}
\hooglig{Cargo} of \remcom{herbs} containing \remcom{whatever} you want.\\\hline
\end{remglyph}

\begin{prons}
\pronsrtwocone&\textbf{カ (KA)}&\textbf{に (ni)}\\\hline
\pronsrthreecone&---&---\\\hline
\end{prons}

\begin{remprons}
\hooglig{The cargo} \comkun{niggled} us to \comon{cuss}.\\\hline
{\warn}\textit{Niggle} means slight but persistent annoyance.\\\hline
\end{remprons}

%{\middel}
%&\mbox{()}\\\hline
\begin{somme}
重荷&heavy + cargo = \textit{heavy burden/load}&おも{\middel}に \mbox{(omo{\middel}ni)}\\\hline
荷重&cargo + heavy = \textit{load; weight}&カ{\middel}ジュウ \mbox{(KA{\middel}JY\={U})}\\\hline
荷物&cargo + thing = \textit{luggage; baggage}&に{\middel}モツ \mbox{(ni{\middel}MOTSU)}\\\hline
出荷&exit + cargo = \textit{shipment; shipping}&シュッ{\middel}カ \mbox{(SHUK{\middel}KA)}\\\hline
入荷&enter + cargo = \textit{goods received}&ニュウ{\middel}カ \mbox{(NY\={U}{\middel}KA)}\\\hline
荷送&cargo + send = \textit{consignment}&に{\middel}おくり \mbox{(ni{\middel}okuri)}\\\hline
着荷&wear + cargo = \textit{arrival of goods}&チャッ{\middel}カ \mbox{(CHAK{\middel}KA)}\\\hline
荷が下りる&cargo + lower = \textit{feeling relieved}&に{\middel}が{\middel}お{\middel}りる \mbox{(ni{\middel}ga{\middel}ori{\middel}ru)}\\\hline
荷が重い&cargo + heavy = \textit{great responsibility}&に{\middel}が{\middel}おも{\middel}い \mbox{(ni{\middel}ga{\middel}omo{\middel}i)}\\\hline
入荷待ち&goods received + wait = \textit{backordering}&ニュウ{\middel}カ{\middel}まち \mbox{(NY\={U}{\middel}KA{\middel}ma{\middel}chi)}\\\hline
\end{somme}

\begin{decomp}{5}
彼&は&荷物&を&下ろした。\\\hline
かれ&は&にモツ&を&おろした\\\hline
he&は&baggage&を&to lower\\\hline
\multicolumn{5}{|r|}{\textit{}}\\\hline
\end{decomp}
\end{multicols}

%\begin{decomp}{}
%\\\hline
%\\\hline
%\\\hline
%\multicolumn{}{|r|}{\textit{}}\\\hline
%\end{decomp}

%\begin{decomp}{4}
%品物&は&昨日&出荷されました。\\\hline
%しなもの&は&きのう&シュッカされました\\\hline
%goods&は&yesterday&shipped\\\hline
%\multicolumn{4}{|r|}{\textit{The goods were sent out yesterday.}}\\\hline
%\end{decomp}
%%--------------------------------------------------------------------
%2--------------------------------------------------------------------
\begin{multicols}{2}
\begin{glyph}{Port (港)}{port}\label{glyph:port}
Port.  Harbour.
\end{glyph}
\gindex{Port}{港}\strindex{12}{港}

\begin{comps}
\comprtwocone&氵 + 共 + 己\\\hline
\comprthreecone&Water + Together (p. \pageref{glyph:together}) + Oneself\\\hline
\end{comps}

\begin{remglyph}
\hooglig{The port} and \remcom{water} \remcom{together} is \remcom{self}-explanatory.\\\hline
\end{remglyph}

\begin{prons}
\pronsrtwocone&\textbf{コウ (K\={O})}&\textbf{みなと (minato)}\\\hline
\pronsrthreecone&---&---\\\hline
\end{prons}

\begin{remprons}
\hooglig{The port} \comon{caused} \comkun{Minato} to flee.\\\hline
{\warn}\textit{Minato} is the father of Naruto in the アニメ.\\\hline
\end{remprons}

%{\middel}
%&\mbox{()}\\\hline
\begin{somme}
港&\textit{port; harbor}&みなと \mbox{(minato)}\\\hline
入港&enter + port = \textit{entering port; arrival}&ニュウ{\middel}コウ \mbox{(NY\={U}{\middel}K\={O})}\\\hline
空港&empty (sky) + port = \textit{airport}&クウ{\middel}コウ \mbox{(K\={U}{\middel}K\={O})}\\\hline
開港&open + port = \textit{opening a port}&カイ{\middel}コウ \mbox{(KAI{\middel}K\={O})}\\\hline
出港&exit + port = \textit{departure from a port}&シュッ{\middel}コウ \mbox{(SHUK{\middel}K\={O})}\\\hline
商業港&commerce + port = \textit{commercial port}&ショウ{\middel}ギョウ{\middel}コウ \mbox{(SH\={O}{\middel}GY\={O}{\middel}K\={O})}\\\hline
港町&port + town = \textit{port city}&みなと{\middel}まち \mbox{(minato{\middel}machi)}\\\hline
港口&port + mouth = \textit{harbour entrance}&コウ{\middel}コウ \mbox{(K\={O}{\middel}K\={O})}\\\hline
港図&port + diagram = \textit{map of a harbour}&コウ{\middel}ズ \mbox{(K\={O}{\middel}ZU)}\\\hline
\end{somme}

\begin{decomp}{4}
今&成田空港&に&います。\\\hline
コン&なりたクウコウ&に&います\\\hline
now&Narita Airport&に&to be\\\hline
\multicolumn{4}{|r|}{\textit{I'm at Narita Airport right now.}}\\\hline
\end{decomp}
\end{multicols}
\pagebreak
%%--------------------------------------------------------------------
%3--------------------------------------------------------------------

\begin{multicols}{2}
\begin{glyph}{Hit (打)}{hit}\label{glyph:hit}
Hit.  Strike.
\end{glyph}
\gindex{Hit}{打}\strindex{5}{打}

\begin{comps}
\comprtwocone&扌 + 丁\\\hline
\comprthreecone&Hand + Block (p. \pageref{glyph:block})\\\hline
\end{comps}

\begin{remglyph}
\hooglig{Hitting} \remcom{hands} in the \remcom{block}.\\\hline
\end{remglyph}

\begin{prons}
\pronsrtwocone&\textbf{ダ (DA)}&\textbf{う (u)}\\\hline
\rye{3}{う is the more common reading.  This is another 漢字 that will often incorporate its accompanying ひらがな when forming nouns.}
\pronsrthreecone&---&---\\\hline
\end{prons}

\begin{remprons}
\hooglig{Hit} the \comon{dud} with a lot of \comkun{oomph}.\\\hline
\end{remprons}

%{\middel}
%&\mbox{()}\\\hline
\begin{somme}
\splitter{打つ}{tr}&\textit{hit; strike}&う{\middel}つ \mbox{(u{\middel}tsu)}\\\hline
\splitter{打ち切る}{tr}&hit + cut = \textit{stop doing sth; break off}&う{\middel}ち{\middel}き{\middel}る \mbox{(u{\middel}chi{\middel}ki{\middel}ru)}\\\hline
打者&hit + individual = \textit{batter; hitter (in baseball)}&ダ{\middel}シャ \mbox{(DA{\middel}SHA)}\\\hline
打ち明ける&hit + bright = \textit{to confide; to open one's heart}&う{\middel}ち{\middel}あ{\middel}ける \mbox{(u{\middel}chi{\middel}a{\middel}keru)}\\\hline
打力&hit + strength = \textit{hitting power}&ダ{\middel}リョク \mbox{(DA{\middel}RYOKU)}\\\hline
頭打ち&head + hit = \textit{reaching a peak/limit}&あたま{\middel}う{\middel}ち \mbox{(atama{\middel}u{\middel}chi)}\\\hline
仕打ち&serve + hit = \textit{poor treatment; bad behaviour}&し{\middel}う{\middel}ち \mbox{(shi{\middel}u{\middel}chi)}\\\hline
早打ち&early + hit = \textit{fast typing/beating}&はや{\middel}う{\middel}ち \mbox{(haya{\middel}u{\middel}chi)}\\\hline
\end{somme}

\begin{decomp}{5}
雨&が&窓&を&打った。\\\hline
あめ&が&まど&を&うった\\\hline
rain&が&window&を&to strike\\\hline
\multicolumn{5}{|r|}{\textit{The rain whipped against the window.}}\\\hline
\end{decomp}
\end{multicols}

%\begin{decomp}{6}
%両国&は&友好&関係&を&打ち切った。\\\hline
%リョウコク&は&ユウコウ&カンケイ&を&うちきった\\\hline
%both countries&は&friendship&relation&を&to stop\\\hline
%\multicolumn{6}{|r|}{\textit{The countries terminated friendly relations.}}\\\hline
%\end{decomp}
%%--------------------------------------------------------------------
%4--------------------------------------------------------------------
\begin{multicols}{2}
\begin{glyph}{Row (列)}{row}\label{glyph:row}
Row (as in a row of objects).
\end{glyph}
\gindex{Row}{列}\strindex{6}{列}

\begin{comps}
\comprtwocone&歹 + 刂\\\hline
\comprthreecone&Death + Knife\\\hline
\end{comps}

\begin{remglyph}
\hooglig{Rows} of \remcom{death and decay} left behind by the \remcom{sword}.\\\hline
\end{remglyph}

\begin{prons}
\pronsrtwocone&\textbf{レツ (RETSU)}&---\\\hline
\pronsrthreecone&---&---\\\hline
\end{prons}

\begin{remprons}
\hooglig{The row} \comon{Retsu} obliterated.\\\hline
{\warn}\textit{Retsu} was the first Kenpachi in the ``Bleach'' アニメ.\\\hline
\end{remprons}

%{\middel}
%&\mbox{()}\\\hline
\begin{somme}
列島&row + island = \textit{island chaino}&レッ{\middel}トウ \mbox{(RET{\middel}T\={O})}\\\hline
列車&row + car = \textit{train}&レッ{\middel}シャ \mbox{(RES{\middel}SHA)}\\\hline
系列&system + row = \textit{series; corporate group}&ケイ{\middel}レツ \mbox{(KEI{\middel}RETSU)}\\\hline
行列&go + row = \textit{line; queue}&ギョウ{\middel}レツ \mbox{(GY\={O}{\middel}RETSU)}\\\hline
最前列&foremost + row = \textit{front row}&サイ{\middel}ゼン{\middel}レツ \mbox{(SAI{\middel}ZEN{\middel}RETSU)}\\\hline
列強&row + strong = \textit{major world powers}&レッ{\middel}キョウ \mbox{(REK{\middel}KY\={O})}\\\hline
系列化&series + change = \textit{putting in order}&ケイ{\middel}レツ{\middel}カ \mbox{(KEI{\middel}RETSU{\middel}KA)}\\\hline
時系列&time + series = \textit{chronological order}&ジ{\middel}ケイ{\middel}レツ \mbox{(JI{\middel}KEI{\middel}RETSU)}\\\hline
\end{somme}

\begin{decomp}{5}
列車&は&六時&に&出る。\\\hline
レッシャ&は&ロクジ&に&でる\\\hline
train&は&6 o'clock&に&to come out\\\hline
\multicolumn{5}{|r|}{\textit{The train starts at six.}}\\\hline
\end{decomp}
\end{multicols}

\pagebreak
%%--------------------------------------------------------------------
%5--------------------------------------------------------------------

\begin{multicols}{2}
\begin{glyph}{Round (丸)}{round}\label{glyph:round}
Round.\\\hline
By extension are the ideas of \textit{completion} and \textit{wholeness}, etc.
\end{glyph}
\gindex{Round}{丸}\strindex{3}{丸}

\begin{comps}
\comprtwocone&九 + 丶\\\hline
\comprthreecone&Nine (p. \pageref{glyph:nine}) + Dot\\\hline
\end{comps}

\begin{remglyph}
\hooglig{Round} \remcom{9} of eating \remcom{Jelly Beans}.\\\hline
\end{remglyph}

\begin{prons}
\pronsrtwocone&---&\textbf{まる (maru)}\\\hline
\pronsrthreecone&ガン (GAN)&---\\\hline
\end{prons}

\begin{remprons}
\hooglig{Round} them up with \ucomon{guns} and take their \comkun{amarula}.\\\hline
\end{remprons}

%{\middel}
%&\mbox{()}\\\hline
\begin{somme}
丸い&\textit{round}&まる{\middel}い \mbox{(maru{\middel}i)}\\\hline
丸める&\textit{to make round}&まる{\middel}める \mbox{(maru{\middel}meru)}\\\hline
日の丸&sun + round = \textit{The Japanese (rising sun) flag}&ひ{\middel}の{\middel}まる \mbox{(hi{\middel}no{\middel}maru)}\\\hline
丸薬&round + medicine = \textit{pill}&ガン{\middel}ヤク \mbox{(GAN{\middel}YAKU)}\\\hline
一丸&one + round = \textit{lump}&イチ{\middel}ガン \mbox{(ICHI{\middel}GAN)}\\\hline
真ん丸い&true + round = \textit{perfectly round}&ま{\middel}ん{\middel}まる{\middel}い \mbox{(ma{\middel}n{\middel}maru{\middel}i)}\\\hline
本丸&main + round = \textit{core; centre; focus}&ホン{\middel}まる \mbox{(HON{\middel}maru)}\\\hline
丸出し&round + exit = \textit{exposing fully}&まる{\middel}だ{\middel}し \mbox{(maru{\middel}da{\middel}shi)}\\\hline
丸見え&round + see = \textit{full view; plain sight}&まる{\middel}み{\middel}え \mbox{(maru{\middel}mi{\middel}e)}\\\hline
\end{somme}

\begin{decomp}{5}
猫&は&背中&を&丸める。\\\hline
ねこ&は&せなか&を&まるめる\\\hline
cat&は&back&を&to curl up\\\hline
\multicolumn{5}{|r|}{\textit{Cats arch their backs.}}\\\hline
\end{decomp}

\end{multicols}

%%--------------------------------------------------------------------
%6--------------------------------------------------------------------
\begin{multicols}{2}
\begin{glyph}{Manor (館)}{manor}\label{glyph:manor}
This 漢字 can appear in compounds referring to anything from stately embassies to simple country inns.
\end{glyph}
\gindex{Manor}{館}\strindex{16}{館}

\begin{comps}
\comprtwocone&食 + 官\\\hline
\comprthreecone&Eat/Food + Official (\ref{glyph:official})\\\hline
\end{comps}

\begin{remglyph}
\hooglig{Manor} where all the \remcom{officials} \remcom{eat}.\\\hline
\end{remglyph}

\begin{prons}
\pronsrtwocone&\textbf{カン (KAN)}&---\\\hline
\pronsrthreecone&---&---\\\hline
\end{prons}

\begin{remprons}
\hooglig{The manor} was \comon{cunningly} built.\\\hline
\end{remprons}

%{\middel}
%&\mbox{()}\\\hline
\begin{somme}
旅館&travel + manor = \textit{Japanese-style inn}&リョ{\middel}カン \mbox{(RYO{\middel}KAN)}\\\hline
図書館&books + manor = \textit{library}&ト{\middel}ショ{\middel}カン \mbox{(TO{\middel}SHO{\middel}KAN)}\\\hline
美術館&art + manor = \textit{art museum}&ビ{\middel}ジュツ{\middel}カン \mbox{(BI{\middel}JUTSU{\middel}KAN)}\\\hline
映画館&movie + manor = \textit{cinema}&エイ{\middel}ガ{\middel}カン \mbox{(EI{\middel}GA{\middel}KAN)}\\\hline
会館&meet + manor = \textit{meeting/assembly hall}&カイ{\middel}カン \mbox{(KAI{\middel}KAN)}\\\hline
本館&main + manor = \textit{main building}&ホン{\middel}カン \mbox{(HON{\middel}KAN)}\\\hline
開館&open + manor = \textit{opening of \ldots}&カイ{\middel}カン \mbox{(KAI{\middel}KAN)}\\\hline
公民館&citizen + manor = \textit{public hall}&コウ{\middel}ミン{\middel}カン \mbox{(K\={O}{\middel}MIN{\middel}KAN)}\\\hline
館内&manor + inside = \textit{in the building}&カン{\middel}ナイ \mbox{(KAN{\middel}NAI)}\\\hline
入館料&entry + fee = \textit{admission fee}&ニュウ{\middel}カン{\middel}リョウ \mbox{(NY\={U}{\middel}KAN{\middel}RY\={O})}\\\hline
\end{somme}

\begin{decomp}{4}
明日&図書館&で&ね。\\\hline
あした&トショカン&で&ね\\\hline
tomorrow&library&で&ね\\\hline
\multicolumn{4}{|r|}{\textit{See you tomorrow at the library.}}\\\hline
\end{decomp}
\end{multicols}

%\begin{decomp}{7}
%図書館&は&市&の&中央&に&ある。\\\hline
%トショカン&は&シ&の&チュウオウ&に&ある\\\hline
%library&は&city&の&centre&に&to be/exist\\\hline
%\multicolumn{7}{|r|}{\textit{The library is in the middle of the city.}}\\\hline
%\end{decomp}

\pagebreak
%%--------------------------------------------------------------------
%7--------------------------------------------------------------------

\begin{multicols}{2}
\begin{glyph}{Fly (飛)}{fly}\label{glyph:fly}
Fly.
\end{glyph}
\gindex{Fly}{飛}\strindex{9}{飛}

\begin{comps}
\comprtwocone&飛\\\hline
\comprthreecone&Fly\\\hline
\end{comps}

\begin{remglyph}
Part of the 漢字 radicals (\#{183}).\\\hline
\end{remglyph}

\begin{prons}
\pronsrtwocone&\textbf{ヒ (HI)}&\textbf{と (to)}\\\hline
\rye{3}{The 訓読み can absorb its accompanying ひらがな when forming nouns.}
\pronsrthreecone&---&---\\\hline
\end{prons}

\begin{remprons}
\hooglig{Fy} away to \comon{heal} your \comkun{torrent} of sorrows.\\\hline
\end{remprons}

%{\middel}
%&\mbox{()}\\\hline
\begin{somme}
\splitter{飛ぶ}{intr}&\textit{to fly}&と{\middel}ぶ \mbox{(to{\middel}bu)}\\\hline
\splitter{売り飛ばす}{tr}&sell + fly = \textit{to dispose of \ldots}&う{\middel}り{\middel}と{\middel}ばす \mbox{(u{\middel}ri{\middel}to{\middel}basu)}\\\hline
飛び出す&fly + exit = \textit{fly out; jump out}&と{\middel}び{\middel}だ{\middel}す \mbox{(to{\middel}bi{\middel}da{\middel}su)}\\\hline
飛行機&flight + mechanism = \textit{aeroplane}&ヒ{\middel}コウ{\middel}キ \mbox{(HI{\middel}K\={O}{\middel}KI)}\\\hline
飛行場&flight + place = \textit{airport}&ヒ{\middel}コウ{\middel}ジョウ \mbox{(HI{\middel}K\={O}{\middel}JY\={O})}\\\hline
飛行士&flight + gentleman = \textit{pilot}&ヒ{\middel}コウ{\middel}シ \mbox{(HI{\middel}K\={O}{\middel}SHI)}\\\hline
飛び切り&fly + cut = \textit{best; unequalled}&と{\middel}び{\middel}き{\middel}り \mbox{(to{\middel}bi{\middel}ki{\middel}ri)}\\\hline
高飛車&tall + fly + car = \textit{high-handed; domineering}&たか{\middel}ビ{\middel}シャ \mbox{(taka{\middel}BI{\middel}SHA)}\\\hline
\end{somme}

\begin{decomp}{5}
楽しい&飛行機&の&旅&を！\\\hline
たのしい&ヒコウキ&の&たび&を\\\hline
enjoyable&aeroplane&の&trip&を\\\hline
\multicolumn{5}{|r|}{\textit{Have a nice flight.}}\\\hline
\end{decomp}
\end{multicols}
%%--------------------------------------------------------------------
%8--------------------------------------------------------------------
\begin{multicols}{2}
\begin{glyph}{Number (数)}{number2}\label{glyph:number}
Number.\\\hline
When in the $1^{\mathrm{st}}$ position, this 漢字 will sometimes carry the meaning of \textit{several} of whatever is indicated by the character following it (e.g. years, persons, days, etc.).
\end{glyph}
\gindex{Number}{数}\strindex{13}{数}

\begin{comps}
\comprtwocone&米 + 女 + 攵\\\hline
\comprthreecone&Rice + Woman + Whip\\\hline
\end{comps}

\begin{remglyph}
\hooglig{Numerous} \remcom{rice} \remcom{women} were \remcom{whipped}.\\\hline
\end{remglyph}

\begin{prons}
\pronsrtwocone&\textbf{スウ (S\={U})}&\textbf{かず、かぞ (kazu, kazo)}\\\hline
\rye{3}{スウ is the most common reading, and the one used to express the idea of \textit{several} mentioned above.\\\hline
かず is the general word for \textit{number} and forms only a few compounds.\\\hline
かぞ is a verb stem.}
\pronsrthreecone&ス (SU)&---\\\hline
\end{prons}

\begin{remprons}
\hooglig{The number} of people following \ucomon{suit} by selling \comkun{kazoos} in the \comkun{car zone} led them to being \comon{sued}.\\\hline
{\warn}\textit{Kazoo} is a small musical instrument.\\\hline
\end{remprons}

%{\middel}
%&\mbox{()}\\\hline
\begin{somme}
\splitter{数える}{tr}&\textit{count}&かぞ{\middel}える \mbox{(kazo{\middel}eru)}\\\hline
数&\textit{number}&かず \mbox{(kazu)}\\\hline
数学&number + study = \textit{mathematics}&スウ{\middel}ガク \mbox{(S\={U}{\middel}GAKU)}\\\hline
人数&person + number = \textit{number of persons}&ニン{\middel}ズウ \mbox{(NIN{\middel}Z\={U})}\\\hline
手数&hand + number = \textit{trouble; bother}&て{\middel}スウ \mbox{(te{\middel}S\={U})}\\\hline
無数&nil + number = \textit{countless}&ム{\middel}スウ　\mbox{(MU{\middel}S\={U})}\\\hline
定数&fixed + number = \textit{constant; fate}&テイ{\middel}スウ \mbox{(TEI{\middel}S\={U})}\\\hline
数日&number + sun = \textit{few days}&スウ{\middel}ジツ \mbox{(S\={U}{\middel}JITSU)}\\\hline
\end{somme}

\begin{decomp}{3}
人数&を&数えます。\\\hline
ニンズウ&を&かぞえます\\\hline
number of people&を&to count\\\hline
\multicolumn{3}{|r|}{\textit{I'm counting the number of people.}}\\\hline
\end{decomp}

\end{multicols}

\pagebreak
%%--------------------------------------------------------------------
%9--------------------------------------------------------------------
\begin{multicols}{2}
\begin{glyph}{Learn (習)}{learn}\label{glyph:learn}
Learn.\\\hline
Take care not to confuse this 漢字 with 皆 (\textit{All}, p. \pageref{glyph:all}).
\end{glyph}
\gindex{Learn}{習}\strindex{11}{習}

\begin{comps}
\comprtwocone&羽 + 白\\\hline
\comprthreecone&Feather + White\\\hline
\end{comps}

\begin{remglyph}
\hooglig{Learn} the \remcom{feather's} \remcom{white} hues.\\\hline
\end{remglyph}

\begin{prons}
\pronsrtwocone&\textbf{シュウ (SH\={U})}&\textbf{なら (nara)}\\\hline
\pronsrthreecone&---&---\\\hline
\end{prons}

\begin{remprons}
\hooglig{Learning} the \comkun{Nara}-brand \comon{shoes}.\\\hline
{\warn}\textit{Nara} is a capital and prefecture name in Japan.\\\hline
\end{remprons}

%{\middel}
%&\mbox{()}\\\hline
\begin{somme}
\splitter{習う}{tr}&\textit{learn; take lessons}&なら{\middel}う \mbox{(nara{\middel}u)}\\\hline
自習&self + learn = \textit{self-study}&ジ{\middel}シュウ \mbox{(JI{\middel}SH\={U})}\\\hline
習字&learn + character = \textit{penmanship}&シュウ{\middel}ジ \mbox{(SH\={U}{\middel}JI)}\\\hline
予習&beforehand + learn = \textit{lesson preparation}&ヨ{\middel}シュウ \mbox{(YO{\middel}SH\={U})}\\\hline
実習&actual + learn = \textit{practice in the field}&ジッ{\middel}シュウ \mbox{(JIS{\middel}SH\={U})}\\\hline
風習&wind + learn = \textit{custom}&フウ{\middel}シュウ \mbox{(F\={U}{\middel}SH\={U})}\\\hline
悪習&evil + learn = \textit{bad habit; vice}&アク{\middel}シュウ \mbox{(AKU{\middel}SH\={U})}\\\hline
習合&learn + match = \textit{syncretism}&シュウ{\middel}ゴウ \mbox{(SH\={U}{\middel}G\={O})}\\\hline
\end{somme}

\begin{decomp}{3}
生け花&を&習いたい。\\\hline
いけばな&を&ならいたい\\\hline
flower arrangement&を&to learn\\\hline
\multicolumn{3}{|r|}{\textit{I'd like to learn how to arrange flowers.}}\\\hline
\end{decomp}
\end{multicols}

%\begin{decomp}{4}
%姉さん&を&見習い&なさい。\\\hline
%ねえさん&を&みならい&なさい\\\hline
%sister&を&to follow another's example&to do\\\hline
%\multicolumn{4}{|r|}{\textit{Follow the example of your sister.}}\\\hline
%\end{decomp}
%%--------------------------------------------------------------------
%10-------------------------------------------------------------------
\begin{multicols}{2}
\begin{glyph}{Display (展)}{display}\label{glyph:display}
Display.　　Unfold.  Spread.
\end{glyph}
\gindex{Display}{展}\strindex{10}{展}

\begin{comps}
\comprtwocone&尸 + 共 + 衣\\\hline
\comprthreecone&Corpse/Body + Together + Clothing\\\hline
\end{comps}

\begin{remglyph}
\hooglig{A display} of the \remcom{corpse} \remcom{together} with its \remcom{clothes}.\\\hline
\end{remglyph}

\begin{prons}
\pronsrtwocone&\textbf{テン (TEN)}&---\\\hline
\pronsrthreecone&---&---\\\hline
\end{prons}

\begin{remprons}
\hooglig{The display} was \comon{tense}.\\\hline
\end{remprons}

%{\middel}
%&\mbox{()}\\\hline
\begin{somme}
展開&display + open = \textit{unfolding; development}&テン{\middel}カイ \mbox{(TEN{\middel}KAI)}\\\hline
展示&display + show = \textit{display; exhibition}&テン{\middel}ジ \mbox{(TEN{\middel}JI)}\\\hline
発展&discharge + display = \textit{expansion}&ハッ{\middel}テン \mbox{(HAT{\middel}TEN)}\\\hline
進展&advance + display = \textit{progress}&シン{\middel}テン \mbox{(SHIN{\middel}TEN)}\\\hline
個展&solo + display = \textit{solo exhibition}&コ{\middel}テン \mbox{(KO{\middel}TEN)}\\\hline
日展&sun + display = \textit{JFAE}&ニッ{\middel}テン \mbox{(NIT{\middel}TEN)}\\\hline
新展&parent + display = \textit{classified information}&シン{\middel}テン \mbox{(SHIN{\middel}TEN)}\\\hline
市場展開&marketplace + unfolding = \textit{market expansion}&シ{\middel}ジョウ{\middel}テン{\middel}カイ \mbox{(SHI{\middel}J\={O}{\middel}TEN{\middel}KAI)}\\\hline
\end{somme}

\begin{decomp}{3}
展示会&を&見てきました。\\\hline
テンジカイ&を&みてきました\\\hline
exhibition&を&to see\\\hline
\multicolumn{3}{|r|}{\textit{We have been to see the exhibition.}}\\\hline
\end{decomp}
\end{multicols}

%\begin{decomp}{9}
%死後&彼&の&絵&は&其の&美術館&に&展示された。\\\hline
%シゴ&かれ&の&エ&は&その&ビジュツカン&に&テンジされた\\\hline
%after death&he&の&painting&は&this/the&art museum&に&exhibition\\\hline
%\multicolumn{9}{|r|}{\textit{After his death, his painting were hung in the museum.}}\\\hline
%\end{decomp}

\pagebreak
%%--------------------------------------------------------------------
%11-------------------------------------------------------------------

\begin{multicols}{2}
\begin{glyph}{Mask (面)}{mask}\label{glyph:mask}
Mask.  Aspect.  Face.\\\hline
The \textit{surface appearance} of things.
\end{glyph}
\gindex{Mask}{面}\strindex{9}{面}

\begin{comps}
\comprtwocone&面\\\hline
\comprthreecone&Mask\\\hline
\end{comps}

\begin{remglyph}
Part of the 漢字 radicals (\#{176}).\\\hline
\end{remglyph}

\begin{prons}
\pronsrtwocone&\textbf{メン (MEN)}&\textbf{おもて、おも (omote, omo)}\\\hline
\pronsrthreecone&---&つら (tsura)\\\hline
\end{prons}

\begin{remprons}
\hooglig{The masks}, washed with \comkun{omo}, by \comon{men} during the \ucomkun{stew lunch} \comkun{commotion}.\\\hline
\end{remprons}

%{\middel}
%&\mbox{()}\\\hline
\begin{somme}
面&\textit{surface; face}&おもて \mbox{(omote)}\\\hline
面白い&mask + white = \textit{interesting}&おも{\middel}しろ{\middel}い \mbox{(omo{\middel}shiro{\middel}i)}\\\hline
面接&mask + join = \textit{interview}&メン{\middel}セツ \mbox{(MEN{\middel}SETSU)}\\\hline
表面&express + mask = \textit{surface; exterior}&ヒョウ{\middel}メン \mbox{(HY\={O}{\middel}MEN)}\\\hline
側面&side + mask = \textit{side; profile}&ソク{\middel}メン \mbox{(SOKU{\middel}MEN)}\\\hline
直面&straight + mask = \textit{confrontation}&チョク{\middel}メン \mbox{(CHOKU{\middel}MEN)}\\\hline
面目&mask + eye = \textit{face; honour; reputation}&メン{\middel}ボク \mbox{(MEN{\middel}BOKU)}\\\hline
\end{somme}

\begin{irrr}
真面目&honesty + eye = \textit{honest; earnest}&ま{\middel}じ{\middel}め \mbox{(ma{\middel}ji{\middel}me)}\\\hline
\end{irrr}

\begin{decomp}{4}
真面目&に&受け取る&な。\\\hline
まじめ&に&うけとる&な\\\hline
earnest&に&to receive&\textit{prohibition}\\\hline
\multicolumn{4}{|r|}{\textit{Don't take it seriously.}}\\\hline
\end{decomp}

\end{multicols}
%%--------------------------------------------------------------------
%12-------------------------------------------------------------------
\begin{multicols}{2}
\begin{glyph}{Select (選)}{select}\label{glyph:select}
Select.  Choose.
\end{glyph}
\gindex{Select}{選}\strindex{15}{選}

\begin{comps}
\comprtwocone&己 + 共 + 辶\\\hline
\comprthreecone&Oneself + Together (p. \pageref{glyph:together}) + Walk/Road\\\hline
\end{comps}

\begin{remglyph}
\hooglig{Select} \remcom{yourself} who you are \remcom{together} with on your \remcom{path}.\\\hline
\end{remglyph}

\begin{prons}
\pronsrtwocone&\textbf{セン (SEN)}&\textbf{えら (era)}\\\hline
\pronsrthreecone&---&---\\\hline
\end{prons}

\begin{remprons}
\hooglig{Choose}, in the \comon{centre}, the \comkun{era} you want to time travel to.\\\hline
\end{remprons}

%{\middel}
%&\mbox{()}\\\hline
\begin{somme}
\splitter{選ぶ}{tr}&\textit{select; choose}&えら{\middel}ぶ \mbox{(era{\middel}bu)}\\\hline
選手&select + hand = \textit{athlete}&セン{\middel}シュ \mbox{(SEN{\middel}SHU)}\\\hline
特選&special + select = \textit{special selection}&トク{\middel}セン \mbox{(TOKU{\middel}SEN)}\\\hline
選考&select + consider = \textit{selection; screening}&セン{\middel}コウ \mbox{(SEN{\middel}K\={O})}\\\hline
予選&beforehand + select = \textit{qualifying round}&ヨ{\middel}セン \mbox{(YO{\middel}SEN)}\\\hline
落選&fall + selection = \textit{election loss; rejection}&ラク{\middel}セン \mbox{(RAKU{\middel}SEN)}\\\hline
選局&select + office = \textit{channel selection; tuning a radio}&セン{\middel}キョク \mbox{(SEN{\middel}KYOKU)}\\\hline
市長選&mayor + select = \textit{mayoral election}&シ{\middel}チョウ{\middel}セン \mbox{(SHI{\middel}CH\={O}{\middel}SEN)}\\\hline
二者選一&two things + select + one = \textit{choosing between two things}&ニ{\middel}シャ{\middel}セン{\middel}イチ \mbox{(NI{\middel}SHA{\middel}SEN{\middel}ICHI)}\\\hline
\end{somme}

\begin{decomp}{3}
此れ&を&選びます。\\\hline
これ&を&えらびます\\\hline
this&を&to select\\\hline
\multicolumn{3}{|r|}{\textit{This is my choice.}}\\\hline
\end{decomp}

\end{multicols}

\pagebreak
%%--------------------------------------------------------------------
%13-------------------------------------------------------------------
\begin{multicols}{2}
\begin{glyph}{Homeward (帰)}{homeward}\label{glyph:homeward}
Homeward, in the sense of \textit{going back} or \textit{returning}.
\end{glyph}
\gindex{Homeward}{帰}\strindex{10}{帰}

\begin{comps}
\comprtwocone&刂 + 帚\\\hline
\comprthreecone&Knife + Broom\\\hline
\rye{2}{帚 = 彐彑 + 冖 + 巾.  Broom = Pig head/snout + Cover/Crown + Cloth.  \hooglig{The broom} is a \remcom{pig's snout} covered in \remcom{cloth}.}
\end{comps}

\begin{remglyph}
\hooglig{Homeward}, I'll return my \textit{sword} and take up a \textit{broom}.\\\hline
\end{remglyph}

\begin{prons}
\pronsrtwocone&\textbf{キ (KI)}&\textbf{かえ (kae)}\\\hline
\rye{3}{The 訓読み becomes voiced according to the same rules as those for 付 (p. \pageref{glyph:attach}).}
\pronsrthreecone&---&---\\\hline
\end{prons}

\begin{remprons}
\hooglig{Homeward} we go after \comon{killing} \comkun{Kaeru-chan}.\\\hline
{\warn}\textit{Kaeru} is a character in the ``ID:Invaded'' \mbox{アニメ}.\\\hline
\end{remprons}

%{\middel}
%&\mbox{()}\\\hline
\begin{somme}
\splitter{持ち帰り}{tr}&hold + homeward = \textit{to bring back}&も{\middel}ち{\middel}かえ{\middel}り \mbox{(mo{\middel}chi{\middel}kae{\middel}ri)}\\\hline
\splitter{帰る}{intr}&\textit{go back}&かえ{\middel}る \mbox{(kae{\middel}ru)}\\\hline
帰り&\textit{return; coming back}&かえ{\middel}り \mbox{(kae{\middel}ri)}\\\hline
日帰り&sun + homeward = \textit{day trip}&ひ{\middel}がえ{\middel}り \mbox{(hi{\middel}gae{\middel}ri)}\\\hline
帰国&homeward + country = \textit{go back to one's country}&キ{\middel}コク \mbox{(KI{\middel}KOKU)}\\\hline
帰化&homeward + change = \textit{naturalisation}&キ{\middel}カ \mbox{(KI{\middel}KA)}\\\hline
帰休&homeward + rest = \textit{leave; furlough}&キ{\middel}キュウ \mbox{(KI{\middel}KY\={U})}\\\hline
帰結&homeward + bind = \textit{consequence; result}&キ{\middel}ケツ \mbox{(KI{\middel}KETSU)}\\\hline
帰着&homeward + wear = \textit{return; conclusion}&キ{\middel}チャク \mbox{(KI{\middel}CHAKU)}\\\hline
\end{somme}

\begin{decomp}{3}
家&に&帰る。\\\hline
いえ&に&かえる\\\hline
house&に&to return\\\hline
\multicolumn{3}{|r|}{\textit{I leave to go home.}}\\\hline
\end{decomp}

\end{multicols}
%%--------------------------------------------------------------------
%14-------------------------------------------------------------------
\begin{multicols}{2}
\begin{glyph}{Section (部)}{section}\label{glyph:section}
Section.  Department.  Category.
\end{glyph}
\gindex{Section}{部}\strindex{11}{部}

\begin{comps}
\comprtwocone&立 + 口 + 阝\\\hline
\comprthreecone&Stand + Mouth + Village/Town/City\\\hline
\end{comps}

\begin{remglyph}
\hooglig{The section} was \remcom{erected} in the \remcom{vampire} \remcom{town}.\\\hline
\end{remglyph}

\begin{prons}
\pronsrtwocone&\textbf{ブ (BU)}&---\\\hline
\pronsrthreecone&---&---\\\hline
\end{prons}

\begin{remprons}
\hooglig{A section} on \comon{bullying}.\\\hline
\end{remprons}

%{\middel}
%&\mbox{()}\\\hline
\begin{somme}
全部&complete + section = \textit{all; the whole}&ゼン{\middel}ブ \mbox{(ZEN{\middel}BU)}\\\hline
上部&upper + section = \textit{upper section/part}&ジョウ{\middel}ブ \mbox{(J\={O}{\middel}BU)}\\\hline
部分&section + section = \textit{part; portion}&ブ{\middel}ブン \mbox{(BU{\middel}BUN)}\\\hline
本部&main + section = \textit{headquarters; head office}&ホン{\middel}ブ \mbox{(HON{\middel}BU)}\\\hline
学部&learn + section = \textit{university department}&ガク{\middel}ブ \mbox{(GAKU{\middel}BU)}\\\hline
部品&section + goods = \textit{parts; components}&ブ{\middel}ヒン \mbox{(BU{\middel}HIN)}\\\hline
部首&section + neck = \textit{kanji radical}&ブ{\middel}シュ \mbox{(BU{\middel}SHU)}\\\hline
\end{somme}

\begin{irrr}
部屋&section + roof = \textit{room (of a house \ldots)}&へ{\middel}や \mbox{(he{\middel}ya)}\\\hline
\end{irrr}

\begin{decomp}{3}
部屋&を&片付けて。\\\hline
へや&を&かたづけて\\\hline
room&を&to tidy up\\\hline
\multicolumn{3}{|r|}{\textit{Clean up the room.}}\\\hline
\end{decomp}

\end{multicols}
\pagebreak
%%--------------------------------------------------------------------
%15-------------------------------------------------------------------

\begin{multicols}{2}
\begin{glyph}{Topic (題)}{topic}\label{glyph:topic}
Topic.  Theme.
\end{glyph}
\gindex{Topic}{題}\strindex{18}{題}

\begin{comps}
\comprtwocone&是 + 頁\\\hline
\comprthreecone&Approve (p. \pageref{glyph:approve}) + Big shell/Page/Face\\\hline
\end{comps}

\begin{remglyph}
\hooglig{Topic} of \remcom{approval} include \remcom{face-recognition}.\\\hline
\end{remglyph}

\begin{prons}
\pronsrtwocone&\textbf{ダイ (DAI)}&---\\\hline
\pronsrthreecone&---&---\\\hline
\end{prons}

\begin{remprons}
\hooglig{The topic} of \comon{dying} is unsettling.\\\hline
\end{remprons}

%{\middel}
%&\mbox{()}\\\hline
\begin{somme}
問題&question + topic = \textit{problem; issue}&モン{\middel}ダイ \mbox{(MON{\middel}DAI)}\\\hline
宿題&lodging + topic = \textit{homework}&シュク{\middel}ダイ \mbox{(SHUKU{\middel}DAI)}\\\hline
主題&primary + topic = \textit{theme; subject}&シュ{\middel}ダイ \mbox{(SHU{\middel}DAI)}\\\hline
話題&speak + topic = \textit{topic; subject}&ワ{\middel}ダイ \mbox{(WA{\middel}DAI)}\\\hline
題名&topic + name = \textit{title; caption}&ダイ{\middel}メイ \mbox{(DAI{\middel}MEI)}\\\hline
出題&exit + topic = \textit{setting a question/theme}&シュツ{\middel}ダイ \mbox{(SHUTSU{\middel}DAI)}\\\hline
社会問題&public + problem = \textit{social problem}&シャ{\middel}カイ{\middel}モン{\middel}ダイ \mbox{(SHA{\middel}KAI{\middel}MON{\middel}DAI)}\\\hline
難題&difficult + topic = \textit{dchallenge}&ナン{\middel}ダイ \mbox{(NAN{\middel}DAI)}\\\hline
別問題&special + issue = \textit{different thing}&ベツ{\middel}モン{\middel}ダイ \mbox{(BETSU{\middel}MON{\middel}DAI)}\\\hline
\end{somme}

\begin{decomp}{3}
問題&だ&なあ。\\\hline
モンダイ&だ&なあ\\\hline
problem&be/is&なあ\\\hline
\multicolumn{3}{|r|}{\textit{It's a problem.}}\\\hline
\end{decomp}

\end{multicols}

%%--------------------------------------------------------------------
%16-------------------------------------------------------------------
\begin{multicols}{2}
\begin{glyph}{Period (期)}{period}\label{glyph:period}
Period (of time).
\end{glyph}
\gindex{Period}{期}\strindex{12}{期}

\begin{comps}
\comprtwocone&其 + 月\\\hline
\comprthreecone&{\warn}That + Moon/Month\\\hline
\end{comps}

\begin{remglyph}
\hooglig{Period} of \remcom{that} \remcom{moon}.\\\hline
\end{remglyph}

\begin{prons}
\pronsrtwocone&\textbf{キ (KI)}&---\\\hline
\pronsrthreecone&ゴ (GO)&---\\\hline
\end{prons}

\begin{remprons}
\hooglig{Period} of \ucomon{gobbling} what you have \comon{killed}.\\\hline
\end{remprons}

%{\middel}
%&\mbox{()}\\\hline
\begin{somme}
期間&period + interval = \textit{period of time}&キ{\middel}カン \mbox{(KI{\middel}KAN)}\\\hline
期待&period + wait = \textit{expectation}&キ{\middel}タイ \mbox{(KI{\middel}TAI)}\\\hline
定期&fixed + period = \textit{fixed period}&テイ{\middel}キ \mbox{(TEI{\middel}KI)}\\\hline
学期&study + period = \textit{school term; semester}&ガッ{\middel}キ \mbox{(GAK{\middel}KI)}\\\hline
周期&around + period = \textit{cycle; period}&シュウ{\middel}キ \mbox{(SH\={U}{\middel}KI)}\\\hline
期日&period + day = \textit{fixed date}&キ{\middel}ジツ\mbox{(KI{\middel}JITSU)}\\\hline
末期&end + period = \textit{end-stage}&マッ{\middel}キ \mbox{(MAK{\middel}KI)}\\\hline
短期&short + period = \textit{short-term}&タン{\middel}キ \mbox{(TAN{\middel}KI)}\\\hline
\notyet{長期}&long + period = \textit{long-term}&チョウ{\middel}キ \mbox{(CH\={O}{\middel}KI)}\\\hline
\end{somme}

\begin{decomp}{5}
スキー&の&時期&は&過ぎた。\\\hline
スキー&の&ジキ&は&すぎた\\\hline
skiing&の&season&は&to pass by\\\hline
\multicolumn{5}{|r|}{\textit{The time for skiing has gone by.}}\\\hline
\end{decomp}

\end{multicols}

%\begin{decomp}{8}
%春休み&は&どれ&くらい&の&期間&です&か。\\\hline
%はるやすみ&は&どれ&くらい&の&キカン&です&か\\\hline
%spring holiday&は&どれ&approximately&の&period of time&be/is&か\\\hline
%\multicolumn{8}{|r|}{\textit{How long is your spring vacation?}}\\\hline
%\end{decomp}

\pagebreak
%%--------------------------------------------------------------------
%17-------------------------------------------------------------------

\begin{multicols}{2}
\begin{glyph}{Grant (授)}{grant}\label{glyph:grant}
To grant.  Impart.
\end{glyph}
\gindex{Grant}{授}\strindex{11}{授}

\begin{comps}
\comprtwocone&扌 + 受\\\hline
\comprthreecone&Hand + Receive (p. \pageref{glyph:receive})\\\hline
\end{comps}

\begin{remglyph}
\hooglig{Grant} your \remcom{actions} to be \remcom{received}.\\\hline
\end{remglyph}

\begin{prons}
\pronsrtwocone&\textbf{ジュ (JU)}&---\\\hline
\pronsrthreecone&---&さず (sazu)\\\hline
\end{prons}

\begin{remprons}
\hooglig{Grants} were \comon{judiciously} given at the \ucomkun{same zoo}.\\\hline
\end{remprons}

%{\middel}
%&\mbox{()}\\\hline
\begin{somme}
授業&grant + business = \textit{lesson; instruction}&ジュ{\middel}ギョウ \mbox{(JU{\middel}GY\={O})}\\\hline
授与&grant + confer = \textit{awarding}&ジュ{\middel}ヨ \mbox{(JU{\middel}YO)}\\\hline
教授&teach + grant = \textit{teaching; a professor}&キョウ{\middel}ジュ \mbox{(KY\={O}{\middel}JU)}\\\hline
授与&grant + confer = \textit{awarding}&ジュ{\middel}ヨ \mbox{(JU{\middel}YO)}\\\hline
天授&heaven + grant = \textit{natural gifts}&テン{\middel}ジュ \mbox{(TEN{\middel}JU)}\\\hline
個人教授&personal + teaching = \textit{private instruction}&コ{\middel}ジン{\middel}キョウ{\middel}ジュ \mbox{(KO{\middel}JIN{\middel}KY\={O}{\middel}JU)}\\\hline
直接教授法&direct + teaching method = \textit{direct method}&チョク{\middel}セツ{\middel}キョウ{\middel}ジュ{\middel}ホウ (CHOKU{\middel}SETSU{\middel}KY\={O}{\middel}JU{\middel}H\={O})\\\hline
\end{somme}

\begin{decomp}{4}
授業&に&遅れる&よ。\\\hline
ジュギョウ&に&おくれる&よ\\\hline
lesson&に&to be late&よ\\\hline
\multicolumn{4}{|r|}{\textit{We'll be late for class.}}\\\hline
\end{decomp}
\end{multicols}

%\begin{decomp}{6}
%授業中&に&話&を&する&な。\\\hline
%ジュギョウチュウ&に&はなし&を&する&な\\\hline
%while in class&に&talk&を&to do&\textit{prohibition}\\\hline
%\multicolumn{6}{|r|}{\textit{Don't speak in the middle of a lesson.}}\\\hline
%\end{decomp}
%%--------------------------------------------------------------------
%18-------------------------------------------------------------------
\begin{multicols}{2}
\begin{glyph}{Delivery (便)}{delivery}\label{glyph:delivery}
This character has an array of meanings, from \textit{postal services}, \textit{convenience} and opportunity to \textit{bodily excretion functions}.  Underlying everything, is a sense of \textit{delivering things} or of \textit{things being delivered} in some way.
\end{glyph}
\gindex{Delivery}{便}\strindex{9}{便}

\begin{comps}
\comprtwocone&亻 + 更\\\hline
\comprthreecone&Man/Person + Anew (p. \pageref{glyph:anew})\\\hline
\end{comps}

\begin{remglyph}
\hooglig{Delivery} makes a \remcom{person} feel \remcom{anew}.\\\hline
\end{remglyph}

\begin{prons}
\pronsrtwocone&\textbf{ベン、ビン (BEN, BIN)}&---\\\hline
\rye{3}{ベン relates to both \textit{convenience} and \textit{bodily functions}.\\\hline
ビン relates to \textit{mail} and \textit{opportunity}.}
\pronsrthreecone&---&たよ (tayo)\\\hline
\end{prons}

\begin{remprons}
\hooglig{Delivered} from \comon{binge} drinking; I \comon{bended} in the direction of \ucomkun{Thai Yoga}.\\\hline
\end{remprons}

%{\middel}
%&\mbox{()}\\\hline
\begin{somme}
便利&delivery + profit = \textit{convenient; handy}&ベン{\middel}リ \mbox{(BEN{\middel}RI)}\\\hline
便所&delivery + location = \textit{toilet; lavatory}&ベン{\middel}ジョ \mbox{(BEN{\middel}JO)}\\\hline
至急便&urgent + delivery = \textit{express (mail)}&シ{\middel}キュウ{\middel}ビン \mbox{(SHI{\middel}KY\={U}{\middel}BIN)}\\\hline
不便&not + convenient = \textit{inconvenience}&フ{\middel}ベン \mbox{(FU{\middel}BEN)}\\\hline
小便&small + delivery = \textit{urine}&ショウ{\middel}ベン \mbox{(SH\={O}{\middel}BEN)}\\\hline
大便&large + delivery = \textit{faeces}&ダイ{\middel}ベン \mbox{(DAI{\middel}BEN)}\\\hline
定期便&regular + delivery = \textit{regular service}&テイ{\middel}キ{\middel}ビン \mbox{(TEI{\middel}KI{\middel}BIN)}\\\hline
用便&use + delivery = \textit{relieving oneself}&ヨウ{\middel}ベン \mbox{(Y\={O}{\middel}BEN)}\\\hline
便意&delivery + nature = \textit{call of nature}&ベン{\middel}イ \mbox{(BEN{\middel}I)}\\\hline
軽便&light + delivery = \textit{simplicity}&ケイ{\middel}ベン \mbox{(KEI{\middel}BEN)}\\\hline
\end{somme}

\begin{decomp}{6}
此れ&は&大変&便利な&物&です。\\\hline
これ&は&タイヘン&ベンリな&もの&です\\\hline
this&は&very&convenient&thing&be/is\\\hline
\multicolumn{6}{|r|}{\textit{This is very useful.}}\\\hline
\end{decomp}

\end{multicols}
\pagebreak
%%--------------------------------------------------------------------
%19-------------------------------------------------------------------
\begin{multicols}{2}
\begin{glyph}{Annual (歳)}{annual}\label{glyph:annual}
Annual.  Year.  Age.\\\hline
When used as a suffix, this 漢字 translates as \textit{{\ldots} years old}.
\end{glyph}
\gindex{Annual}{歳}\strindex{13}{歳}

\begin{comps}
\comprtwocone&止 + 戌 + 小\\\hline
\comprthreecone&Stop + Sign of the Dog + Small\\\hline
\rye{2}{戌 = 戈 + 一.  Sign of the Dog = Spear/Halberd + One.  \hooglig{The Sign of the Dog} is close to the \remcom{one} \remcom{halberd}.}
\end{comps}

\begin{remglyph}
\hooglig{Annually} we \remcom{stop} to gaze at the \remcom{Sign of the Dog}, though \remcom{small}.\\\hline

\end{remglyph}

\begin{prons}
\pronsrtwocone&\textbf{サイ (SAI)}&---\\\hline
\pronsrthreecone&セイ (SEI)&---\\\hline
\end{prons}

\begin{remprons}
\hooglig{Annual} \ucomon{sailing} \comon{psyches} me up.\\\hline
\end{remprons}

%{\middel}
%&\mbox{()}\\\hline
\begin{somme}
何歳&what + annual = \textit{how old?}&なん{\middel}サイ \mbox{(nan{\middel}SAI)}\\\hline
万歳&$10,000$ + annual = \textit{Hurrah!}&バン{\middel}ザイ \mbox{(BAN{\middel}ZAI)}\\\hline
万々歳&$10,000$ $\left( \times 2 \right)$ + annual = \textit{matter for great celebration}&バン{\middel}バン{\middel}ザイ \mbox{(BAN{\middel}BAN{\middel}ZAI)}\\\hline
五歳&five + annual = \textit{five years old}&ゴ{\middel}サイ \mbox{(GO{\middel}SAI)}\\\hline
歳出&annual + exit = \textit{annual expenditure}&サイ{\middel}シュツ \mbox{(SAI{\middel}SHUTSU)}\\\hline
歳入&annual + enter = \textit{annual revenue}&サイ{\middel}ニュウ \mbox{(SAI{\middel}NY\={U})}\\\hline
歳月&annual + moon = \textit{time; years}&サイ{\middel}ゲツ \mbox{(SAI{\middel}GETSU)}\\\hline
%歳末&annual + end = \textit{year end}&サイ{\middel}マツ \mbox{(SAI{\middel}MATSU)}\\\hline
\end{somme}

\begin{irrr}
二十歳&twenty + annual = \textit{twenty years old}&はたち \mbox{(hatachi)}\\\hline
\end{irrr}

\begin{decomp}{4}
ところで、&何歳&です&か。\\\hline
ところで&なんサイ&です&か\\\hline
incidentally&how old&be/is&か\\\hline
\multicolumn{4}{|r|}{\textit{By the way, how old are you?}}\\\hline
\end{decomp}
\end{multicols}
%%--------------------------------------------------------------------
%20-------------------------------------------------------------------
\begin{multicols}{2}
\begin{glyph}{Exclusive (専)}{exclusive}\label{glyph:exclusive}
Exclusive.
\end{glyph}
\gindex{Exclusive}{専}\strindex{9}{専}

\begin{comps}
\comprtwocone&一 + 由 + 寸\\\hline
\comprthreecone&One + Cause + Inch/Tiny\\\hline
\end{comps}

\begin{remglyph}
\hooglig{Exclusively} for \remcom{one} \remcom{tiny} \remcom{cause}.\\\hline
\end{remglyph}

\begin{prons}
\pronsrtwocone&\textbf{セン (SEN)}&---\\\hline
\pronsrthreecone&---&もっぱ (moppa)\\\hline
\end{prons}

\begin{remprons}
\hooglig{Exclusive} \ucomkun{moppers} available at any \comon{centre}.\\\hline
\end{remprons}

%{\middel}
%&\mbox{()}\\\hline
\begin{somme}
専門&exclusive + gate = \textit{speciality}&セン{\middel}モン \mbox{(SEN{\middel}MON)}\\\hline
専用&exclusive + utilize = \textit{exclusive use; private use}&セン{\middel}ヨウ \mbox{(SEN{\middel}Y\={O})}\\\hline
専有&exclusive + have = \textit{exclusive possession}&セン{\middel}ユウ \mbox{(SEN{\middel}Y\={U})}\\\hline
専制&exclusive + regulation = \textit{despotism; autocracy}&セン{\middel}セイ \mbox{(SEN{\middel}SEI)}\\\hline
専門学校&speciality + school = \textit{technical school}&セン{\middel}モン{\middel}ガッ{\middel}コウ \mbox{(SEN{\middel}MON{\middel}GAK{\middel}K\={O})}\\\hline
専門的&speciality + target = \textit{technical}&セン{\middel}モン{\middel}テキ \mbox{(SEN{\middel}MON{\middel}TEKI)}\\\hline
一意専心&one though + undivided attention = \textit{single-mindedly; wholeheartedly}&イチ{\middel}イ{\middel}セン{\middel}シン \mbox{(ICHI{\middel}I{\middel}SEN{\middel}SHIN)}\\\hline
%専決&exclusive + decide = \textit{deciding; acting on one's own}&セン{\middel}ケツ \mbox{(SEN{\middel}KETSU)}\\\hline
\end{somme}

\begin{decomp}{5}
専攻&は&何&です&か。\\\hline
センコウ&は&なに&です&か\\\hline
major subj.&は&what&be/is&か\\\hline
\multicolumn{5}{|r|}{\textit{What is your major?}}\\\hline
\end{decomp}
\end{multicols}

%\begin{decomp}{6}
%彼&は&天文学&の&専門家&だ。\\\hline
%かれ&は&テンモンガク&の&センモンカ&だ\\\hline
%he&は&astronomy&の&specialist&be/is\\\hline
%\multicolumn{6}{|r|}{\textit{He is an expert in astronomy.}}\\\hline
%\end{decomp}
\pagebreak
%%--------------------------------------------------------------------
%21-------------------------------------------------------------------

\begin{multicols}{2}
\begin{glyph}{Orient (向)}{orient}\label{glyph:orient}
To orient.  To look on.  The overall sense is of things either \textit{orienting themselves}, or \textit{being oriented} in a literal or figurative direction.\\\hline
When used as a suffix, the meaning becomes \textit{intended for} or \textit{suited to}.
\end{glyph}
\gindex{Orient}{向}\strindex{6}{向}

\begin{comps}
\comprtwocone&ノ + 冂 + 口\\\hline
\comprthreecone&NO + To Enclose/Wilderness + Mouth\\\hline
\end{comps}

\begin{remglyph}
\hooglig{Orient} the \remcom{Noh} mask \remcom{to enclose} your \remcom{mouth}.\\\hline
\end{remglyph}

\begin{prons}
\pronsrtwocone&\textbf{コウ (K\={O})}&\textbf{む (mu)}\\\hline
\rye{3}{む usually acts as a verb stem.}
\pronsrthreecone&---&---\\\hline
\end{prons}

\begin{remprons}
\hooglig{Orient} yourself to the \comkun{moon} regarding your \comon{cause}.\\\hline
\end{remprons}

%{\middel}
%&\mbox{()}\\\hline
\begin{somme}
\splitter{向ける}{tr}&to turn towards \ldots&む{\middel}ける \mbox{(mu{\middel}keru)}\\\hline
\splitter{向かう}{intr}&\textit{face on; be opposite}&む{\middel}かう \mbox{(mu{\middel}kau)}\\\hline
向こう&\textit{other side}&む{\middel}こう \mbox{(mu{\middel}k\={o})}\\\hline
夏向き&summer + orient = \textit{for summer}&なつ{\middel}む{\middel}き \mbox{(natsu{\middel}mu{\middel}ki)}\\\hline
向学心&orient + study + heart = \textit{love of learning}&コウ{\middel}ガク{\middel}シン \mbox{(K\={O}{\middel}GAKU{\middel}SHIN)}\\\hline
意向&mind + orient = \textit{inclination}&イ{\middel}コウ \mbox{(I{\middel}K\={O})}\\\hline
動向&move + orient = \textit{trend; tendency; attitude}&ドウ{\middel}コウ \mbox{(D\={O}{\middel}K\={O})}\\\hline
転向&roll + orient = \textit{conversion; switch}&テン{\middel}コウ \mbox{(TEN{\middel}K\={O})}\\\hline
風向&wind + orient = \textit{wind direction}&フウ{\middel}コウ \mbox{(F\={U}{\middel}K\={O})}\\\hline
\end{somme}

\begin{decomp}{5}
二人&は&面&と&向かった。\\\hline
ニニン&は&メン&と&むかった\\\hline
two people&は&face&と&to face on\\\hline
\multicolumn{5}{|r|}{\textit{They confronted each other.}}\\\hline
\end{decomp}

\end{multicols}
%%--------------------------------------------------------------------
%22-------------------------------------------------------------------
\begin{multicols}{2}
\begin{glyph}{Train (練)}{train}\label{glyph:train}
To train.  Knead.
\end{glyph}
\gindex{Train}{練}\strindex{14}{練}

\begin{comps}
\comprtwocone&糸 + 東\\\hline
\comprthreecone&Thread + East (p. \pageref{glyph:east})\\\hline
\end{comps}

\begin{remglyph}
\hooglig{Train} your \remcom{threading} in an \remcom{Eastern} way.\\\hline
\end{remglyph}

\begin{prons}
\pronsrtwocone&\textbf{レン (REN)}&\textbf{ね (ne)}\\\hline
\pronsrthreecone&---&---\\\hline
\end{prons}

\begin{remprons}
\hooglig{Training} \comkun{necromancy} while \comon{renting} a house in the Netherworld.\\\hline
\end{remprons}

%{\middel}
%&\mbox{()}\\\hline
\begin{somme}
\splitter{練る}{tr}&\textit{to train; knead}&ね{\middel}る \mbox{(ne{\middel}ru)}\\\hline
練習&train + learn = \textit{practice; exercise}&レン{\middel}シュウ \mbox{(REN{\middel}SH\={U})}\\\hline
訓練&instruction + train = \textit{training; discipline}&クン{\middel}レン \mbox{(KUN{\middel}REN)}\\\hline
未練&not yet + train = \textit{lingering attachment; regret}&ミ{\middel}レン \mbox{(MI{\middel}REN)}\\\hline
試練&try + train = \textit{test; trial; ordeal}&シ{\middel}レン \mbox{(SHI{\middel}REN)}\\\hline
練達&train + achieve = \textit{expertise; dexterity}&レン{\middel}タツ \mbox{(REN{\middel}TATSU)}\\\hline
朝練&morning + train = \textit{morning training}&あさ{\middel}レン \mbox{(asa{\middel}REN)}\\\hline
軍事教練&military affairs + drill = \textit{military training}&グン{\middel}ジ{\middel}キョウ{\middel}レン \mbox{(GUN{\middel}JI{\middel}KY\={O}{\middel}REN)}\\\hline
訓練所&training + locate = \textit{training school}&クン{\middel}レン{\middel}ジョ \mbox{(KUN{\middel}REN{\middel}JO)}\\\hline
\end{somme}

\begin{decomp}{5}
是非&練習&を&続け&為さい。\\\hline
ゼヒ&レンシュウ&を&つづけ&なさい\\\hline
without fail&practice&を&to continue&to do\\\hline
\multicolumn{5}{|r|}{\textit{Do keep practising!}}\\\hline
\end{decomp}

\end{multicols}

\pagebreak
%%--------------------------------------------------------------------
%23-------------------------------------------------------------------

\begin{multicols}{2}
\begin{glyph}{Raise (育)}{raise}\label{glyph:raise}
To raise.  To bring up (as in children).
\end{glyph}
\gindex{Raise}{育}\strindex{8}{育}

\begin{comps}
\comprtwocone&亠 + 厶 + 月\\\hline
\comprthreecone&Lid/Top + Self/Private + Meat\\\hline
\end{comps}

\begin{remglyph}
\hooglig{Raise} the \remcom{lid} in \remcom{private} as if it's your own \remcom{flesh}.\\\hline
\end{remglyph}

\begin{prons}
\pronsrtwocone&\textbf{イク (IKU)}&\textbf{そだ (soda)}\\\hline
\pronsrthreecone&---&---\\\hline
\end{prons}

\begin{remprons}
\hooglig{Raise} your children by \comon{eke-ing} out the \comkun{soda} usage.\\\hline
\end{remprons}

%{\middel}
%&\mbox{()}\\\hline
\begin{somme}
\splitter{育てる}{tr}&\textit{raise; bring up}&そだ{\middel}てる \mbox{(soda{\middel}teru)}\\\hline
教育&teach + raise = \textit{education}&キョウ{\middel}イク \mbox{(KY\={O}{\middel}IKU)}\\\hline
教育者&education + individual = \textit{educator}&キョウ{\middel}イク{\middel}シャ \mbox{(KY\={O}{\middel}IKU{\middel}SHA)}\\\hline
体育館&physical education + manor = \textit{gymnasium}&タイ{\middel}イク{\middel}カン \mbox{(TAI{\middel}IKU{\middel}KAN)}\\\hline
育成&raise + become = \textit{rearing}&イク{\middel}セイ \mbox{(IKU{\middel}SEI)}\\\hline
生育&life + raise = \textit{birth and growth}&セイ{\middel}イク \mbox{(SEI{\middel}IKU)}\\\hline
発育&discharge + raise = \textit{physical growth}&ハツ{\middel}イク \mbox{(HATSU{\middel}IKU)}\\\hline
育休&raise + rest = \textit{childcare leave}&イク{\middel}キュウ \mbox{(IKU{\middel}KY\={U})}\\\hline
\end{somme}

\begin{decomp}{5}
教育&は&家庭&に&始まる。\\\hline
キョウイク&は&カテイ&に&はじまる\\\hline
training&は&home&に&to begin\\\hline
\multicolumn{5}{|r|}{\textit{Education starts at home.}}\\\hline
\end{decomp}
\end{multicols}

%\begin{decomp}{5}
%彼ら&は&体育館&に&集まります。\\\hline
%かれか&は&タイイクカン&に&あつまります\\\hline
%they&は&gymnasium&に&to gather\\\hline
%\multicolumn{5}{|r|}{\textit{They gathered in the gym.}}\\\hline
%\end{decomp}
%%--------------------------------------------------------------------
%24-------------------------------------------------------------------
\begin{multicols}{2}
\begin{glyph}{Consider (考)}{consider}\label{glyph:consider}
Consider.  Think about.
\end{glyph}
\gindex{Consider}{考}\strindex{6}{考}

\begin{comps}
\comprtwocone&耂 + 丂\\\hline
\comprthreecone&Old + Reach\\\hline
\end{comps}

\begin{remglyph}
\hooglig{Consider} the \remcom{old} ways to \remcom{reach} your potential.\\\hline
\end{remglyph}

\begin{prons}
\pronsrtwocone&\textbf{コウ (K\={O})}&\textbf{かんが (kanga)}\\\hline
\pronsrthreecone&---&---\\\hline
\end{prons}

\begin{remprons}
\hooglig{Consider} the \comon{cause} of the \comkun{kangaroos'} assault on us.\\\hline
\end{remprons}

%{\middel}
%&\mbox{()}\\\hline
\begin{somme}
\splitter{考える}{tr}&\textit{consider; think about}&かんが{\middel}える \mbox{(kanga{\middel}eru)}\\\hline
考え&\textit{thinking; thought}&かんが{\middel}え \mbox{(kanga{\middel}e)}\\\hline
思考&think + consider = \textit{thinking; consideration}&シ{\middel}コウ \mbox{(SHI{\middel}K\={O})}\\\hline
考え方&consider + direction = \textit{way of thinking}&かんが{\middel}え{\middel}かた \mbox{(kanga{\middel}e{\middel}kata)}\\\hline
考古学&consider + old + study = \textit{archeology}&コウ{\middel}コ{\middel}ガク \mbox{(K\={O}{\middel}KO{\middel}GAKU)}\\\hline
選考&select + consider = \textit{selection; screening}&セン{\middel}コウ \mbox{(SEN{\middel}K\={O})}\\\hline
考え直す&consider + straight = \textit{to reconsider}&かんが{\middel}え{\middel}なお{\middel}す \mbox{(kanga{\middel}e{\middel}nao{\middel}su)}\\\hline
失考&lose + consider = \textit{misunderstanding}&シッ{\middel}コウ \mbox{(SHIK{\middel}K\={O})}\\\hline
千思万考&1,000 + think + 10,000 + consider = \textit{deep meditation}&センシバンコウ \mbox{(SEN{\middel}SHI{\middel}BAN{\middel}K\={O})}\\\hline
\end{somme}

\begin{decomp}{2}
考えて&下さい。\\\hline
かんがえて&ください\\\hline
to think about&to kindly do for one\\\hline
\multicolumn{2}{|r|}{\textit{Think about it.}}\\\hline
\end{decomp}

\end{multicols}
\pagebreak
%%--------------------------------------------------------------------